% !PRINTED 274
\weberpage{274}
Für $n=3$ erhält man aus (13):
\weq{20}{
x^2y^2 + x'^2y'^2 = 1,
}
eine Form der Modulargleichung, die von \gesperrt{Legendre} herrührt.

Daraus, indem man für $x'$, $y'$ aus (19) die Werte setzt,
\[
x^8 + y^8 = 4x^2y^2 - 6x^4y^4 + 4x^6y^6,
\]
oder
\[
(x^4 - y^4)^2 = 4(xy - x^3y^3)^2
\]
und indem man hieraus die Wurzel zieht und das Vorzeichen
durch $q=0$ bestimmt, findet man:
\weq{21}{
x^4 - y^4 = 2xy - 2x^3y^3, \qquad (n=3),
}
was, abgesehen von dem dort nicht näher erklärten Vorzeichen
von $\sqrt[4]{\lambda}$, mit \S 10, (12) übereinstimmt.

Für $n=5$ folgt aus der zweiten Gleichung (18):
\[
(x^2 - y^2)^3 = 4xy\,x'^4y'^4,
\]
und daraus durch Quadrieren
\[
(x^2 - y^2)^6 = 16x^2y^2(1 - x^8 - y^8 + x^8y^8),
\]
was leicht in die Form gebracht wird
\[
(x^2 - y^2)^2\bigl[(x^2 - y^2)^2 + 8x^2y^2\bigr]^2 = 16x^2y^2(1 - x^4y^4)^2.
\]
Zieht man hieraus die Wurzel, so folgt, wie oben:
\weq{22}{
x^6 - y^6 - 4xy(1 - x^4y^4) + 5x^2y^2(x^2 - y^2) = 0 \quad (n=5).
}
Für $n=7$ erhält man, ohne Wurzelziehen, aus (8):
\weq{23}{
xy + x'y' = 1,
}
und daraus:
\weq{24}{
\begin{aligned}
x^8 + y^8 - 8xy(1 + x^6y^6) &+ 28x^2y^2(1 + x^4y^4)\\
- 56x^3y^3(1 + x^2y^2) &+ 70x^4y^4 = 0 \qquad (n=7).
\end{aligned}
}

\begin{center}
\S\ 76.\quad \gesperrt{Zusammengesetzte Transformationsgrade.}
\end{center}

Ist der Transformationsgrad eine zusammengesetzte Zahl, so
kann man noch einfachere Transformationsgleichungen aufstellen
als die, die man auf dem Wege des vorigen Paragraphen gewinnt.

Wir führen diese Betrachtungen hier nur in den einfachsten
Fällen durch.

Der ungerade Transformationsgrad $n$ sei in zwei Faktoren
zerlegt
\weq{1}{
n = n'n'',
}
die zueinander relativ prim sind.
